\documentclass[11pt,a4paper]{article}
\usepackage[utf8]{inputenc}
\usepackage[T1]{fontenc}
\usepackage{amsmath,amssymb,amsthm}
\usepackage{graphicx}
\usepackage{booktabs}
\usepackage{hyperref}
\usepackage[margin=1in]{geometry}
\usepackage{parskip}
\usepackage{natbib}

\title{Comprehensive Framework for Quantum-Secured Terrestrial Point-to-Point Radio Backhauling: Metamaterials, Quantum Illumination, and Radiative Cooling Architectures}
\author{}
\date{}

\begin{document}
\maketitle

\begin{abstract}
The trajectory of telecommunications infrastructure is currently undergoing a profound paradigm shift. As the industry accelerates toward the deployment of sixth-generation (6G) communication systems, the underlying architecture must support unprecedented performance metrics, targeting peak data rates between 100 Gbps and 1 Tbps, while simultaneously reducing latency to 0.1 milliseconds. While highly resilient optical fiber networks form the backbone of global connectivity, the economic and geographic realities of planetary topology dictate that a pervasive digital divide cannot be eradicated through wired infrastructure alone. Remote rural terrains, dense urban microcells, and maritime environments fundamentally require the deployment of wireless alternatives. Consequently, terrestrial Point-to-Point (P2P) radio backhauling has emerged as an indispensable pillar of modern network architecture, designed to establish multi-gigabit wireless connections between distributed small cells, macrocells, and the central core network without the prohibitive capital expenditure of physical cable trenching.
\end{abstract}

\section{Introduction}

The trajectory of telecommunications infrastructure is currently undergoing a profound paradigm shift. As the industry accelerates toward the deployment of sixth-generation (6G) communication systems, the underlying architecture must support unprecedented performance metrics, targeting peak data rates between 100 Gbps and 1 Tbps, while simultaneously reducing latency to 0.1 milliseconds.\cite{ref1} While highly resilient optical fiber networks form the backbone of global connectivity, the economic and geographic realities of planetary topology dictate that a pervasive digital divide cannot be eradicated through wired infrastructure alone.\cite{ref3} Remote rural terrains, dense urban microcells, and maritime environments fundamentally require the deployment of wireless alternatives.\cite{ref3} Consequently, terrestrial Point-to-Point (P2P) radio backhauling has emerged as an indispensable pillar of modern network architecture, designed to establish multi-gigabit wireless connections between distributed small cells, macrocells, and the central core network without the prohibitive capital expenditure of physical cable trenching.\cite{ref3}

Concurrently, the rapid maturation of quantum computing threatens to render classical cryptographic encryption obsolete, necessitating the integration of quantum communication protocols---specifically the distribution of continuous-variable (CV) quantum states---directly into the backhaul physical layer.\cite{ref6} Historically, long-distance quantum communication research has heavily favored Satellite Communications (SATCOM) to leverage the near-vacuum of space, despite its severe restrictions involving atmospheric Mie scattering, strict launch load constraints, and orbital mechanics.\cite{ref6} Transitioning these quantum protocols to a macroscopic, terrestrial P2P radio network introduces a substantially more hostile propagation environment. Terrestrial microwave links are subjected to intense urban multipath scattering, atmospheric gas absorption, severe precipitation attenuation, and the overwhelming thermal noise penalty inherent to the 300 Kelvin (K) ambient Earth environment.\cite{ref2}

To decouple the scaling of secure P2P terrestrial backhaul from these severe environmental and thermodynamic overloads, advanced network engineering mandates the integration of ``Quantum Application-Specific Integrated Circuits'' (Quantum ASICs) utilizing programmable holographic metasurfaces.\cite{ref6} By shifting from static physical antennas to dynamically reconfigurable, software-defined quantum arrays developed through rigorous Engineering-as-Code (EaC) pipelines, operators can manipulate photonic and phononic states at the sub-wavelength scale.\cite{ref6} This exhaustive research report analyzes the computational physics, signal processing models, and hardware ecosystems required to deploy resilient, quantum-enhanced P2P radio backhaul networks, bridging the immense gap between deep-cryogenic quantum state generation and chaotic terrestrial deployment.

\section{Environmental Impediments in Terrestrial P2P Radio Backhaul}

Deploying continuous-variable quantum states across terrestrial P2P links demands a rigorous understanding of the environmental factors that actively degrade radio frequency (RF) propagation. Classical 5G and 6G backhaul networks rely on aggressive spectrum expansion, migrating from sub-6 GHz regimes into microwave (6--42 GHz), millimeter-wave (mmWave, 30--300 GHz), and sub-terahertz (sub-THz, 100 GHz to 1 THz) frequency bands to secure wide-channel bandwidths.\cite{ref4} While these higher frequencies unlock vast data capacities, their compressed wavelengths render the physical layer extraordinarily susceptible to dynamic atmospheric and topological channel impairments.\cite{ref14}

\subsection{Atmospheric Absorption and Rain Attenuation Models}

Unlike fiber optic cables, which isolate optical photons within total internal reflection cores, unguided terrestrial radio signals must navigate the molecular makeup of the Earth's atmosphere.\cite{ref10} In terrestrial microwave and mmWave links, atmospheric gases---predominantly oxygen and water vapor---possess discrete molecular resonance frequencies that strongly absorb RF energy.\cite{ref12} This interaction converts the electromagnetic signal into molecular kinetic energy, resulting in a distance-dependent exponential decay of signal strength.\cite{ref10}

Furthermore, as carrier frequencies scale above 10 GHz, rain attenuation (commonly referred to as rain fade) emerges as the dominant fading mechanism.\cite{ref2} Under heavy precipitation conditions, the physical diameter of falling rain droplets begins to approximate the operational wavelength of the transmitted signal.\cite{ref12} This geometric parity induces severe Rayleigh and Mie scattering alongside direct absorption, stripping immense amounts of energy from the communication channel.\cite{ref6} Accurate link budget analyses for these terrestrial backhaul links rely on specialized deterministic and empirical path loss models.\cite{ref16} Planners must utilize localized ITU-R rain attenuation time-series models, taking into account specific regional thunderstorm ratios, liquid water content, and complex permittivity to calculate the requisite fade margins and outage probabilities necessary to maintain carrier-grade link availability.\cite{ref16}

\subsection{Multipath Scattering: Rayleigh and Rician Fading Distributions}

Beyond molecular and atmospheric absorption, the physical topology of the terrestrial landscape severely distorts the spatiotemporal integrity of the signal. The radio link between a transmitting small cell base station and a receiving backhaul node is rarely an unobstructed free-space channel.\cite{ref19} Instead, it is continuously obstructed by dense urban architecture, terrain features, and dynamic vehicular traffic.\cite{ref19} This structural clutter causes the transmitted electromagnetic wave to reflect, refract, and diffract, spawning multiple delayed copies of the original signal.\cite{ref10} These multipath components arrive at the receiver antenna shifted in amplitude, frequency (due to Doppler spread), and phase.\cite{ref23}

The constructive and destructive interference of these delayed electromagnetic waves is mathematically modeled through statistical fading distributions:

\textbf{Rayleigh Fading:} This statistical model accurately characterizes heavily built-up urban macro-environments where there is absolutely no dominant Line-of-Sight (LOS) propagation path between the transmitter and receiver.\cite{ref25} Because the received signal is entirely comprised of scattered, independent multipath components reflected off surrounding structures, the central limit theorem dictates that the channel impulse response is best modeled as a complex Gaussian process.\cite{ref25} In the absence of a dominant vector, this process exhibits a zero mean, with a phase that is evenly distributed between $-\pi$ and $\pi$ radians.\cite{ref25} The resulting envelope of the received signal fluctuates violently, inducing severe amplitude drops (deep fades) that catastrophically degrade the performance of classical coherent receivers relying on phase stability.\cite{ref19}

\textbf{Rician Fading:} Conversely, when a direct LOS path is maintained in conjunction with the scattered multipath components, the channel statistics shift to a Rician distribution.\cite{ref25} In a Rician fading channel, the complex-valued channel coefficient is divided into two distinct parts: a dominant, deterministic primary signal (the LOS path) and a scattered, Rayleigh-distributed component.\cite{ref28} The superposition of these elements yields a Gaussian distribution with a prominent, non-zero mean.\cite{ref27} The severity of the fading is dictated by the Rician $K$-factor, which represents the ratio of the dominant signal power to the variance of the multipath scatter.\cite{ref28} As dynamic obstacles obscure the LOS path, the dominant component decays, and the Rician distribution seamlessly degrades back into a severe Rayleigh profile.\cite{ref29}

\begin{table}[htbp]
\centering
\caption{Fading model comparison for terrestrial receivers.}
\begin{tabular}{@{}llll@{}}
\toprule
Fading Model & Dominant Propagation & Statistical Mean & Impact on Terrestrial Receiver \\
\midrule
Rayleigh Fading & Non-Line-of-Sight (NLOS), pure scatter & Zero Mean & Severe signal attenuation, deep interference nulls, high outage probability \\
Rician Fading & Line-of-Sight (LOS) + scattered paths & Non-Zero Mean & Moderate fading; requires dynamic channel estimation to maintain phase \\
\bottomrule
\end{tabular}
\end{table}

In classical telecommunications, mitigation techniques such as spatial diversity, adaptive modulation, and Massive Multiple-Input Multiple-Output (MIMO) array processing are employed to combat these multipath effects.\cite{ref30} However, when attempting to transmit highly fragile continuous-variable quantum entangled states across these terrestrial P2P links, classical mitigation falls drastically short. The severe Signal-to-Noise Ratio (SNR) degradation caused by Rayleigh scattering rapidly breaks the delicate entanglement properties required for secure quantum backhaul, necessitating entirely new quantum-domain detection protocols.\cite{ref33}

\section{The Thermodynamic Bottleneck and Solid-State Quantum ASICs}

Resolving the extreme signal attenuation caused by terrestrial fading requires massive, highly directive phased-array antennas capable of manipulating the local density of photonic states at the sub-wavelength scale.\cite{ref6} However, integrating advanced quantum components into these arrays introduces a profound thermodynamic bottleneck. The fundamental physical challenge in deploying a quantum-secured microwave backhaul link lies in the exceptionally small energy of microwave photons, which is governed by the Planck-Einstein relation ($E = h\nu$).\cite{ref6} Because the energy scales of standard backhaul frequencies (e.g., 6 GHz to 42 GHz) are minute, the underlying control layers responsible for generating and routing these states must operate well below 100 millikelvin (mK).\cite{ref5}

This deep-cryogenic requirement is necessary to prevent thermal occlusion.\cite{ref6} Thermal occlusion occurs when the ambient thermal energy of the operational environment, dictated by the Boltzmann constant ($k_B$), approaches or exceeds the energy of the quantum state being manipulated.\cite{ref6} When this energy threshold is crossed, the coherent quantum backhaul signal is effectively masked by a dense wash of thermally excited environmental photons, leading to rapid decoherence and the irreversible loss of secure cryptographic information.\cite{ref6}

\subsection{The Coaxial Wiring Crisis vs.\ 28nm FD-SOI Cryo-CMOS}

Operating a massive 1,000-element quantum metasurface array for highly directive P2P beamforming utilizing conventional control architectures is physically impossible due to the ``coaxial wiring crisis''.\cite{ref6} Driving modern superconducting circuits currently mandates a brute-force connectivity approach, utilizing numerous physical coaxial interconnects routed from room-temperature digital-to-analog converters down to the deepest millikelvin stages of a dilution refrigerator.\cite{ref6} This static wiring paradigm dissipates active power and conducts passive heat, creating an insurmountable thermodynamic burden.\cite{ref6} Even when utilizing highly specialized low-thermal-conductivity alloys, a single UT-085-SS-SS coaxial cable introduces nearly 1 milliwatt (mW) of heat load to the 4 K stage of a cryogenic system.\cite{ref6} When this heat load is multiplied by the thousands of interconnects required for a dense macro-scale antenna, it instantly overwhelms the micro-watt cooling capacities available at the lowest cryogenic stages.\cite{ref6} Furthermore, traditional optical phase-shifting materials, such as nematic liquid crystals, freeze completely at these temperatures, losing all orientational mobility.\cite{ref6}

To resolve this wiring crisis and enable terrestrial P2P metasurface architectures, proposed designs demand the use of heterogeneously integrated Cryo-CMOS controllers co-located directly with the quantum hardware at 10 mK.\cite{ref6} Standard bulk CMOS transistors fail catastrophically at cryogenic temperatures due to dopant freeze-out; near absolute zero, thermal energy is insufficient to ionize dopant atoms, causing charge carriers to become permanently bound and turning the semiconductor substrate into an insulator.\cite{ref6}

To overcome this materials science limitation, theoretical cryogenic quantum controllers rely on 28-nanometer Fully Depleted Silicon-On-Insulator (28nm FD-SOI) technology.\cite{ref6} In an FD-SOI architecture, the active transistor channel is formed by an ultra-thin layer of silicon that is entirely undoped.\cite{ref6} Because there are no dopants to freeze out, the transistor remains highly conductive and perfectly operational even at 100 mK.\cite{ref6} These controllers function as massive digital multiplexers, replacing thousands of analog physical cables with a highly simplified digital input signal, such as a 4-wire Serial Peripheral Interface (SPI), routed from a warmer room-temperature stage.\cite{ref6}

Upon reaching the deep-cryogenic Cryo-CMOS chip, an on-board digital finite-state machine interprets the SPI instructions and configures an array of analog circuit blocks known as Charge-Lock Fast-Gate (CLFG) cells.\cite{ref6} At 10 mK, the off-state leakage current of an FD-SOI transistor is practically non-existent.\cite{ref6} This unique physical property allows the CLFG cells to trap exact amounts of charge on floating capacitors to generate stable, static bias voltages without drawing continuous current.\cite{ref6} Simulated performance profiles demonstrate an unprecedented active power dissipation of merely 18 nanowatts (nW) per cell when generating 100 millivolt (mV) control pulses.\cite{ref6} Consequently, a fully scaled array of 1,000 individual control channels would dissipate only 18 microwatts ($\mu$W), remaining safely below the thermal failure threshold of the metasurface's cryogenic housing.\cite{ref6}

\begin{table}[htbp]
\centering
\caption{Control architecture comparison: traditional coaxial routing vs.\ 28nm FD-SOI Cryo-CMOS.}
\begin{tabular}{@{}lll@{}}
\toprule
Control Architecture Aspect & Traditional Coaxial Routing & Proposed 28nm FD-SOI Cryo-CMOS \\
\midrule
Physical Interconnects & 1 dedicated coax line per meta-atom element & 4-wire digital Serial Peripheral Interface (SPI) \\
Heat Load to Cryogenic Stages & $\sim$1 milliwatt (mW) per line & 18 nanowatts (nW) per CLFG cell \\
Signal Generation Mechanism & Room temperature DACs + heavy attenuation & On-chip Charge-Lock Fast-Gate capacitors \\
Scalability Constraint & Hard limit reached via cooling power exhaustion & Highly scalable, minimal dissipation footprint \\
\bottomrule
\end{tabular}
\end{table}

\section{Quantum Design Automation (QDA) and Engineering-as-Code}

The transition of the quantum-enabled terrestrial backhaul antenna from a theoretical physics construct into a physically deployable architecture cannot rely on manual CAD layout processes.\cite{ref6} To achieve the combinatorial scale required for high-bandwidth, fault-tolerant communication arrays, the telecommunications industry is aggressively shifting toward Quantum Design Automation (QDA) workflows.\cite{ref6} QDA provides a systematic, automated computational framework for the design, simulation, and verification of massive superconducting circuits, acting as the necessary quantum-centric analog to classical Electronic Design Automation (EDA).\cite{ref6}

This paradigm shift is heavily underpinned by the ``Engineering-as-Code'' (EaC) and ``Materials-as-Code'' philosophies.\cite{ref6} EaC dictates that every aspect of the physical infrastructure---from the geometric layout of thousands of superconducting microwave traces to the cloud environments hosting the multi-physics simulations---must be explicitly expressed as human-readable, declarative code.\cite{ref6} By treating physical hardware definitions identically to application software, quantum RF engineers can subject antenna designs to rigorous version control, peer review, and automated integration testing, preventing configuration drift across highly complex development pipelines.\cite{ref6}

\subsection{Automated Layout Generation via Python Frameworks}

At the foundation of the QDA stack are programmatic layout generators such as Qiskit Metal and KQCircuits.\cite{ref6} These Python-based software development kits (SDKs) allow researchers to bypass manual drawing by instantiating complex superconducting elements through high-level application programming interfaces (APIs).\cite{ref6} By defining parameters in Python dictionaries---such as specifying a RouteMeander with a total length of 6 mm and specific pin inputs connecting an OpenToGround termination---the software automatically calculates the required geometries.\cite{ref6} The automation algorithms inherently respect microwave design constraints, autonomously managing the complex filleting of trace corners and the alignment of lead straights to maintain a strict 50-ohm characteristic impedance across the entire routing topology.\cite{ref6}

Once the Python scripts execute, these frameworks compile the programmatic definitions into Graphic Data System (GDSII) files, the industry-standard format for integrated circuit photomask production.\cite{ref6} Auxiliary Python packages further automate the post-processing of these layouts, programmatically converting patterns for positive resist exposure, distributing microscopic holes for under-etching processes, generating alignment markers, and exporting structures directly into 3D rendering environments.\cite{ref6}

\subsection{GitOps for Superconducting Circuit Verification}

Treating macroscopic backhaul hardware as code necessitates the implementation of modern software delivery mechanisms, specifically GitOps.\cite{ref6} GitOps applies continuous integration and continuous deployment (CI/CD) practices directly to the hardware infrastructure, establishing a Git repository as the single, immutable source of truth for the metasurface's declarative state.\cite{ref6} The foundational principles of GitOps---declarative configuration, versioned immutability, pull-based deployments, and continuous reconciliation---are critical for mathematically managing the combinatorial complexity of 1,000-element phased arrays.\cite{ref6}

In a theoretical QDA GitOps pipeline, an engineer modifies a Python script to adjust the capacitive coupling of a specific subset of antenna elements. Upon committing this code to the repository, a CI runner is automatically triggered.\cite{ref6} Rather than immediately transmitting the design to a semiconductor foundry, the CI pipeline executes the script to generate the new GDSII layout within a secure, ephemeral container.\cite{ref6} The pipeline then autonomously initiates Design Rule Checks (DRC) and Layout-Versus-Schematic (LVS) verifications.\cite{ref6} Utilizing logical parameters defined directly within the foundry's Process Design Kit (PDK), the system ensures that the generated geometry does not violate minimum trace widths, cryogenic layer spacing rules, or required exclusion zones.\cite{ref6}

If the automated verification detects a physical design violation, the pipeline immediately fails, blocking the pull request and providing diagnostic feedback to the engineer without requiring manual oversight.\cite{ref6} If the checks pass, a GitOps reconciliation agent automatically promotes the verified GDSII artifact to a container registry or cloud storage bucket, making it available for subsequent electromagnetic simulation stages.\cite{ref6} This fully automated promotion flow mathematically guarantees that only structurally sound metamaterial designs progress through the development pipeline, eliminating flawed architectures before physical fabrication is ever attempted.\cite{ref6}

\section{Computational Electrodynamics of Cryogenic Meta-Atoms}

The programmable holographic metasurface designed for terrestrial microwave backhaul requires pure solid-state, dissipationless components that function as flux-tunable meta-atoms.\cite{ref6} These advanced metamaterials must exhibit highly tunable reactive properties---specifically variable inductance or capacitance---driven entirely by external magnetic flux or electrostatic fields generated by the underlying 28nm FD-SOI Cryo-CMOS logic.\cite{ref6} Two specific classes of solid-state meta-atoms satisfy these stringent thermodynamic constraints: radio-frequency Superconducting Quantum Interference Devices (rf-SQUIDs) and piezoelectric lithium niobate (LiNbO$_3$) Bulk Acoustic Wave (BAW) resonators.\cite{ref6}

\subsection{Simulating Quantum Energy Spectra with scqubits}

At the microscopic level, the primary active component enabling dynamic microwave phase shifting across the P2P array is the rf-SQUID.\cite{ref6} Unlike a direct-current (DC) SQUID, which contains two Josephson junctions, an rf-SQUID consists of a single superconducting loop interrupted by exactly one highly nonlinear Josephson junction.\cite{ref6} The operational principle relies heavily on the quantization of magnetic flux and the nonlinear kinetic inductance provided by the junction.\cite{ref6} In a superconductor, the inertia of the Cooper pairs carrying the supercurrent contributes an additional inductance term that is strictly dependent on the gauge-invariant phase difference across the junction.\cite{ref6} Because this phase difference is dictated by the total external magnetic flux threading the loop, the effective inductance of the rf-SQUID can be radically altered by changing its local magnetic environment.\cite{ref6}

To extract the fundamental energy spectra, wavefunctions, and coherence metrics of these individual meta-atoms, computational physicists heavily utilize \texttt{scqubits}, an open-source Python library dedicated to the numerical analysis of superconducting circuits.\cite{ref6} The Hamiltonian for a generic flux-tunable superconducting loop is programmatically defined within the software as:
\begin{equation}
H = 4E_C n^2 - E_J \cos(\phi) + \frac{1}{2} E_L(\phi - \phi_{\text{ext}})^2
\end{equation}
In this computational model, $E_C$ represents the charging energy, $E_J$ is the Josephson energy, $E_L$ is the inductive energy of the loop, $\phi$ is the phase across the junction, $n$ is the conjugate charge operator, and $\phi_{\text{ext}}$ is the external magnetic flux bias.\cite{ref6} By truncating the theoretically infinite-dimensional Hilbert space to a computationally manageable localized basis, scqubits utilizes high-performance SciPy backend eigensolvers to rapidly diagonalize the massive Hamiltonian matrix.\cite{ref6} Automated execution of sweeping parameters allows engineers to algorithmically identify exact magnetic flux bias ``sweet spots.'' At these coordinates, the software confirms that the non-linearity of the kinetic inductance is maximized while the first derivative of the transition frequency with respect to flux approaches zero.\cite{ref6} Operating at these sweet spots actively minimizes the meta-atom's susceptibility to ambient $1/f$ flux noise, ensuring high-fidelity operation without requiring physical trial-and-error in a cryostat.\cite{ref6}

\subsection{Macroscopic 2D London Equation Solvers}

While scqubits effectively resolves the internal quantum dynamics of an isolated meta-atom, designing the macroscopic Quantum ASIC requires scaling the simulation to a massive 1,000-element metasurface array.\cite{ref6} Transitioning from a single component to a macroscopic terrestrial antenna fundamentally alters the computational approach.\cite{ref6} Traditional 3D Finite Element Method (FEM) solvers often struggle with the extreme aspect ratios characteristic of thin-film superconductors.\cite{ref6} Attempting to mesh an ultra-thin superconducting layer within a large spatial bounding box leads to combinatorial mesh explosions, consuming excessive RAM and resulting in intractable computation times.\cite{ref6}

To bypass this computational bottleneck, the Engineering-as-Code roadmap leverages specialized 2D solvers such as SuperScreen, an open-source Python package designed specifically to simulate the magnetic response of quasi-2D superconductors.\cite{ref6} Rather than solving full 3D Maxwell equations, SuperScreen computationally solves the 2D London equation using an exceptionally efficient matrix inversion method.\cite{ref6} The software computes the spatial distribution of Meissner screening currents and trapped magnetic flux given an applied, time-independent magnetic field.\cite{ref6} The foundational physics equation relates the 2D sheet current $\mathbf{K}$ to the local magnetic vector potential $\mathbf{A}$:
\begin{equation}
\mathbf{K}(x,y) = -\frac{1}{\mu_0 \Lambda(x,y)} \left[ \mathbf{A}(x,y) + \frac{\Phi_0}{2\pi} \nabla\theta(x,y) \right]
\end{equation}
In this framework, $\Lambda(x,y)$ represents the spatially inhomogeneous effective magnetic penetration depth, $\Phi_0$ is the magnetic flux quantum, $\mu_0$ is the vacuum permeability, and $\theta$ is the superconducting phase.\cite{ref6} A profound advantage of SuperScreen is its programmatic ability to handle complex, irregular device geometries with spatially varying $\Lambda$ parameters.\cite{ref6} This is essential for modeling realistic rf-SQUID arrays where minor fabrication imperfections or variations in the Ginzburg-Landau critical temperature introduce localized disorder across the metasurface.\cite{ref6} By generating optimized finite element data structures, SuperScreen rapidly calculates the massive self-inductance and mutual-inductance matrices of the entire aperiodic array, bridging the gap between localized quantum kinetics and macroscopic antenna electrodynamics.\cite{ref6}

\subsection{Open Quantum Systems and TLS Decoherence Mitigation}

Complementing the purely inductive response of superconducting rf-SQUIDs, the architecture also relies on advanced acoustic metamaterials. While microwave photons are excellent for fast communication across the backhaul link, their long wavelengths make sub-millimeter spatial localization difficult.\cite{ref6} Phonons (quantized acoustic vibrations) operating at identical gigahertz frequencies possess wavelengths roughly five orders of magnitude smaller than their electromagnetic counterparts, allowing for incredibly dense, highly localized meta-atom arrays within the P2P antenna.\cite{ref6} Piezoelectric lithium niobate (LiNbO$_3$) Bulk Acoustic Wave (BAW) resonators are highly attractive due to their strong electromechanical coupling coefficients.\cite{ref6}

When operated in the high-frequency radio regime and cooled to standard cryogenic temperatures (4 K), acoustic attenuation in the pristine crystal lattice is primarily governed by the Landau-Rumer microscopic model of phonon-phonon scattering.\cite{ref6} Based strictly on these mathematical models, the mechanical quality factor ($Q$) should rise toward infinity as thermal phonon populations drop exponentially near absolute zero.\cite{ref6} However, empirical data replicated in advanced computational models reveals a profound material anomaly: the quality factor violently degrades when the simulation environment is computationally pushed down to 10 mK to match the operational environment of the superconducting qubits.\cite{ref6}

The physical mechanism responsible for this anomaly is resonant absorption by microscopic Two-Level System (TLS) defects.\cite{ref6} At higher temperatures, thermal energy is high enough that microscopic structural impurities---such as dangling bonds or oxygen vacancies within the crystal bulk---are constantly saturated and transitioning, rendering them largely transparent.\cite{ref6} However, as the temperature drops to 10 mK, these defects freeze into their absolute ground states, forming highly coherent Two-Level Systems.\cite{ref6} When the meta-atom is excited by a backhaul transmission, the propagating phonons or microwave photons are resonantly absorbed by these frozen defects, acting as a devastating source of dielectric noise, energy relaxation ($T_1$ decay), and phase decoherence ($T_2$ decay).\cite{ref6}

Understanding and mitigating this noise requires explicit code-based material science modeling the hardware as an open quantum system.\cite{ref6} Because TLS defects represent a massive environmental bath interacting with the primary quantum system, their dynamics cannot be modeled by simple unitary Schrödinger evolution.\cite{ref6} Instead, computational physicists track the time evolution of the system's density matrix $\rho$ using the Lindblad master equation:
\begin{equation}
\dot{\rho} = -\frac{i}{\hbar}[H_{\text{sys}} + H_{\text{bath}} + H_{\text{int}}, \rho] + \sum_k \left( L_k \rho L_k^\dagger - \frac{1}{2}\{L_k^\dagger L_k, \rho\} \right)
\end{equation}
In this programmatic model, $L_k$ are the collapse (jump) operators describing irreversible energy loss and dephasing to the macroscopic environment.\cite{ref6} Open-source libraries like the Quantum Toolbox in Python (QuTiP) provide built-in \texttt{DecoherenceNoise} classes to numerically integrate the master equation over time.\cite{ref6} By executing these models, engineers mathematically validate that actively tuning TLS frequencies out of the spectral vicinity of the signal via applied local DC electric fields can effectively force the parasitic defects into an acoustic vacuum, computationally predicting coherence improvements of up to two orders of magnitude prior to physical fabrication.\cite{ref6}

\section{Machine Learning for Real-Time Holography and Quantum Routing}

Deploying these meticulously designed elements into a functional P2P backhaul node requires highly sophisticated algorithmic frameworks capable of dynamically programming the holographic metasurface to create transient electromagnetic channels that bridge logical data streams to physical transmission antennas.\cite{ref6} Because physical components on a solid-state chip can only interact with their immediate nearest neighbors, bridging non-adjacent nodes requires dynamically calculating the most efficient spatial routing topologies.

Determining the optimal spatial assignment of logical roles to physical qubits to minimize total gate depth and SWAP count is an NP-hard problem.\cite{ref6} To computationally resolve this, the routing challenge is mathematically formulated as a Quadratic Unconstrained Binary Optimization (QUBO) problem.\cite{ref6} Binary decision variables representing specific edge routings are mapped to a rigorous cost function designed to heavily penalize inefficient wiring distances and deep circuit depths that would expose the quantum state to TLS decoherence before transmission completes.\cite{ref6} The optimization is executed using hybrid quantum-classical variational algorithms, such as the Quantum Approximate Optimization Algorithm (QAOA), which navigates the massive multidimensional Hilbert space to extract optimal logical-to-physical mapping vectors.\cite{ref6}

\subsection{Real-Time Phase Generation via PyTorch Deep Neural Networks}

Translating these low-dimensional routing topologies into high-dimensional, continuous electromagnetic phase commands for the 1,000 individual meta-atoms poses the most severe computational bottleneck in the terrestrial architecture.\cite{ref6} Traditional computational electromagnetics rely on solving full-wave Maxwell equations using Finite-Difference Time-Domain (FDTD) methods.\cite{ref6} Calculating complex scattering matrices for an aperiodic array takes hours or days on classical supercomputers---extreme latency that is fundamentally incompatible with the microsecond-scale coherence times ($T_2$) of superconducting quantum states.\cite{ref6}

To achieve real-time dynamic holography and bypass the FDTD bottleneck entirely, the system leverages Deep Neural Network (DNN) surrogate modeling implemented via PyTorch.\cite{ref6} Deployed through robust Machine Learning Operations (MLOps) pipelines, a massive dataset is autonomously generated using standard FDTD solvers running on scalable compute clusters.\cite{ref6} This data correlates specific meta-atom geometric configurations to their resulting S-parameters.\cite{ref6}

A Convolutional Neural Network (CNN) Forward Prediction model is trained on this dataset, rigorously minimizing the Mean Squared Error (MSE) between the FEM-simulated parameters and the network's predictions.\cite{ref6} Once the forward model is validated, an Inverse Design Network utilizes non-linear activation functions like Rectified Linear Units (ReLU) to construct highly complex latent representations of the desired backhaul scattering matrix.\cite{ref6} The final output layer maps this representation to a continuous tensor representing the exact physical phase shifts required for all meta-atoms simultaneously, executing real-time wavefront steering in milliseconds.\cite{ref6}

\subsection{Thermodynamic Validation: The $\pi$ Radian Baseline}

A profound operational insight emerges from the empirical benchmarking of this generative PyTorch model. When synthesizing the massive element phase profile required to steer the microwave beam toward a receiving P2P backhaul node, the neural network autonomously generates phase shifts strictly concentrated in a highly narrow band between 3.033 and 3.284 radians, with a precise mean phase of 3.142 radians ($\pi$).\cite{ref6}

This algorithmic output represents a monumental thermodynamic triumph.\cite{ref6} A uniform $\pi$ phase shift applied across an entire surface effectively acts as a perfect flat inversion mirror.\cite{ref6} By computationally applying only subtle, sub-radian perturbations around this $\pi$-radian baseline, the metasurface leverages slight geometric phase deviations to create the precise constructive and destructive interference patterns required to focus the communication beam.\cite{ref6} Because the AI surrogate restricts the required phase tuning to such a narrow margin, the underlying Cryo-CMOS controllers are not required to inject massive amounts of voltage to drive broad phase sweeps.\cite{ref6} This computational constraint explicitly preserves the ultra-low 18 nanowatt active dissipation limit per cell; without it, physical actuation of broad phase shifts would generate excess Joule heating, instantly crashing the cryostat via thermal occlusion.\cite{ref6}

\section{Quantum Microwave Illumination: Conquering Multipath Clutter}

While the Quantum ASIC provides the internal computational engine, the physical layer of the terrestrial radio link must overcome the severe environmental impediments discussed previously. The severe signal degradation imposed by Rayleigh fading and urban multipath clutter actively destroys the delicate entanglement properties required for secure backhaul.\cite{ref19} Classical mitigation falls short, requiring an entirely new quantum-domain detection framework.\cite{ref33}

To overcome these physical limitations, the quantum-enabled terrestrial backhaul architecture implements protocols derived from Microwave Quantum Illumination (QI).\cite{ref34} Microwave QI represents a fundamental paradigm shift in low-power wireless detection and communication, purposefully designed for environments where traditional coherent radars and classical heterodyne receivers fail due to overwhelming thermal noise and intense multi-bounce scattering.\cite{ref34}

\subsection{Generating Two-Mode Squeezed Vacuum (TMSV) States}

The core physical engine of the quantum illumination protocol relies on the generation of highly correlated photon pairs exhibiting continuous-variable entanglement. This is achieved through the deterministic creation of Two-Mode Squeezed Vacuum (TMSV) states via spontaneous parametric down-conversion.\cite{ref33} In the physical hardware layer, this process utilizes Josephson Parametric Amplifiers (JPAs) situated within the deep-cryogenic stages of the metasurface.\cite{ref34}

The non-degenerate parametric amplification entangles a microwave ``signal'' mode ($a_S$) with an ``idler'' mode ($a_I$), driven by a pump field operating at twice the signal frequency ($2\omega$).\cite{ref33} The quantum Hamiltonian for this nonlinear process, $H_{\text{int}} \propto a_S a_I + a_S^\dagger a_I^\dagger$, illustrates the robust creation of these coupled photon statistics.\cite{ref33} In the terrestrial P2P backhaul execution, the signal photon is transmitted out of the metasurface aperture, navigating the turbulent, multipath-dominated atmospheric channel toward the receiving backhaul node.\cite{ref37} Concurrently, the idler photon is perfectly retained locally within a cryogenic quantum memory at the transmitter.\cite{ref36}

As the signal photon travels through the urban environment, severe attenuation ($\eta \ll 1$) and inescapable interactions with the 300K thermal Earth background completely break the pure quantum entanglement between the signal and idler pairs.\cite{ref33}

\subsection{Phase-Conjugate Receivers and the Chernoff Error Exponent}

Despite the catastrophic loss of pure entanglement in the transit channel, mathematical simulations of non-classical correlation metrics prove that a robust, residual signature---specifically quantum discord---survives the entanglement-breaking noise.\cite{ref6} When the heavily degraded, multipath-scattered return signal is captured by the receiving aperture, it cannot be processed by standard classical thresholding.\cite{ref33} Instead, the architecture employs advanced quantum receiver systems, such as optical parametric amplifier-single-photon avalanche diode (OPA-SPADE) configurations or digital phase-conjugation modules.\cite{ref33}

These receivers execute Joint Positive-Operator-Valued Measures (POVMs), performing a joint measurement of the arriving scattered signal photon against the locally stored idler photon.\cite{ref33} By exploiting the surviving quantum discord, this joint measurement effectively extracts cross-correlations (idler-conditional signal statistics) that reject common-mode noise far better than classical direct imaging.\cite{ref33}

Performance analyses under realistic thermal noise models ($n_{\text{th}} \gg 1$) prove that this protocol substantially boosts the quantum Chernoff error exponent ($\xi_{\text{QI}}$), which dictates the asymptotic error probability.\cite{ref33} This mathematical advantage elevates detection precision to the theoretical Heisenberg limit, empowering the microwave metasurface to reject urban multipath clutter that is up to 20 dB stronger than the true target echoes.\cite{ref6} This capability doubles the effective communication and detection range compared to the best possible classical heterodyne receiver operating at equivalent milliwatt transmit powers, validating QI's necessity in dense 6G deployments.\cite{ref6} Furthermore, optimizing the waveforms with chirped pulses and Hermite-Gaussian modes concentrates entanglement within specific bandwidth-time products, further reducing peak power and boosting effective detection efficiency under multipath conditions.\cite{ref33}

\section{Radiative Sky Cooling and Elevated Temperature Quantum Receivers}

While quantum illumination actively addresses the spatial scattering of the signal, establishing an unconditional quantum backhaul link requires directly neutralizing the immense thermodynamic penalty of the terrestrial environment itself. The fatal flaw of attempting microwave quantum communication through the Earth's atmosphere is ambient thermal noise.\cite{ref6} At standard terrestrial temperatures (approximately 300 K), the environment naturally generates a massive population of thermal microwave photons that instantly drown out fragile continuous-variable states attempting to cross the channel.\cite{ref6} While SATCOM systems point outward toward the 2.7 K vacuum of space to naturally mitigate some of this noise, horizontal P2P terrestrial links do not benefit from this cosmic cold load.\cite{ref6}

\subsection{The 8--13\,\textmu m Atmospheric Transparency Window}

To actively drain this thermal noise out of the terrestrial receiver without relying on massive, power-intensive active refrigeration, advanced thermodynamic architectures integrate ``radiative sky cooling'' directly into the metasurface and antenna housings.\cite{ref43} Radiative cooling utilizes the fundamental principles of thermal radiation; heat flows from a hotter object to a colder sink via electromagnetic waves without requiring an intermediary material medium.\cite{ref41}

The Earth's atmosphere uniquely possesses a primary infrared transmission window between 8 and 13 micrometers ($\mu$m).\cite{ref45} Because this specific spectral window aligns perfectly with the peak of the blackbody radiation curve for objects near 300 K, a tailored material can emit thermal radiation directly through the atmosphere and into the extreme 3 K vacuum of deep space, without the radiation being absorbed or reflected back by atmospheric gases.\cite{ref41}

Using a thermal photonic approach, engineers design integrated solar reflectors and thermal emitters utilizing stacked nanophotonic structures (such as alternating layers of hafnium dioxide, HfO$_2$, and silicon dioxide, SiO$_2$).\cite{ref45} These highly selective emitters reflect up to 97 percent of incident sunlight while emitting strongly within the 8--13 $\mu$m atmospheric window.\cite{ref50} This allows the physical housing of the receiver to passively drop to temperatures substantially below the ambient air, providing continuous, electricity-free cooling power up to 140 W/m$^2$ to pre-cool the system.\cite{ref41}

\subsection{Thermal Photon Suppression via 10 mK Overcoupling}

While passive radiative cooling handles the external structural loads, the internal signal transmission lines rely on dynamic, active cooling protocols to manage the volatile microwave photon statistics. Recent breakthroughs utilize 4-Kelvin transmission lines composed of niobium--titanium (NbTi).\cite{ref44} NbTi is a superconductor specifically selected for its low loss and robust quantum coherence at elevated temperatures (4 K), far above the milli-Kelvin regimes required by standard aluminum materials.\cite{ref44} Despite the transmission line being physically thermalized at 4 K, researchers apply an active ``radiative cooling'' method directly to the microwave channel.\cite{ref44}

In standard thermodynamics, an antenna naturally reaches thermal equilibrium with its operational environment. By computationally overcoupling the microwave communication channel to an ancillary cold load maintained at 10 mK (situated deep within the system's localized dilution refrigerator), the architecture creates an immense thermal gradient.\cite{ref6} The strong external coupling channel rapidly dissipates the chaotic thermal photons from the ``hot'' 4 K environment directly into the 10 mK reservoir.\cite{ref43}

Mathematical models and empirical validation confirm that this radiative overcoupling successfully manipulates the Bose-Einstein statistics of the channel, driving the effective thermal photon occupation number down to an exceptional 0.06.\cite{ref6} This two-orders-of-magnitude reduction makes the 4 K channel operate as if it were in a deep cryogenic regime, fully protected from the 300K terrestrial noise floor.\cite{ref44} To execute a transmission, the system dynamically decouples the channel from the cold load, transferring the quantum state near-instantaneously before the transmission line can physically rethermalize and induce decoherence.\cite{ref44}

Further supplementing this architecture is the deployment of nanobridge kinetic-inductance parametric amplifiers (NKPA) patterned from high critical-temperature ($T_c$) Niobium Nitride (NbN) thin films.\cite{ref43} Unlike traditional aluminum-based Josephson amplifiers that are strictly constrained below 1 K, the NbN NKPA can be ``uplifted'' to warmer cryogenic stages (1.5 K to 4 K).\cite{ref43} When radiatively cooled via an external coupling bath, these amplifiers maintain a high 27 dB gain at 1.5 K while introducing an added noise of merely 1.3 quanta---which represents only a 0.2 quanta increase compared to baseline operations at 0.1 K.\cite{ref43} This interconnected cooling architecture dramatically reduces the power and space constraints of the dilution fridges, permitting unconditionally secure state transfer across terrestrial channels that would otherwise destroy the signal.\cite{ref6}

\section{Continuous-Variable Quantum Teleportation (CVQT) in P2P Topologies}

With Rayleigh multipath clutter successfully suppressed by quantum illumination protocols, and the ambient 300K thermal noise heavily mitigated by radiative sky cooling and cold-load overcoupling, the physical layer of the P2P backhaul is finally primed for secure data transmission. However, to completely bypass the unavoidable physical attenuation caused by atmospheric gas absorption and rain fade over macroscopic distances (often spanning several kilometers), the network cannot rely on direct signal transmission alone; it must utilize Continuous-Variable Quantum Teleportation (CVQT).\cite{ref51}

Unlike discrete-variable teleportation which relies on single photons (which are statistically difficult to generate deterministically and require highly inefficient, expensive single-photon avalanche detectors), CVQT utilizes the continuous degrees of freedom of the electromagnetic field---specifically the amplitude (``position'' $q$) and phase (``momentum'' $p$) quadratures.\cite{ref52}

The CVQT teleportation protocol begins by distributing a maximally entangled resource state between the two P2P backhaul nodes.\cite{ref54} This is typically a TMSV state produced by interacting two infinitely-squeezed vacuum states ($\lvert 0\rangle_q$ and $\lvert 0\rangle_p$) on a highly calibrated 50:50 beam splitter.\cite{ref54} To teleport an unknown coherent quantum state $\lvert\alpha\rangle$ across the link, the transmitting node (Alice) performs a joint Bell state measurement.\cite{ref54} She interferes the target state $\lvert\alpha\rangle$ with her local copy of the entangled resource state using a secondary 50:50 beam splitter, feeding the output directly into a homodyne detection array.\cite{ref54}

This measurement process is inherently entangling; Alice then transmits the classical description of her measurement outcome via a standard, high-bandwidth radio link to the receiving backhaul node (Bob).\cite{ref54} Using this classical data, Bob applies a series of precise conditional phase and amplitude displacements to his retained half of the entangled state, perfectly recovering the original quantum state $\lvert\alpha\rangle$ without the physical carrier itself ever traversing the lossy atmospheric channel.\cite{ref54}

\subsection{Non-Gaussian Operations and NLA Integration}

While standard CVQT protocols rely heavily on Gaussian resource states (such as the standard TMSV), terrestrial networks subjected to intense fading require enhanced resilience. Advanced research demonstrates that executing specific non-Gaussian operations---such as photon subtraction, photon addition, or the implementation of ``quantum scissors''---prior to transmission can drastically alter the Wigner function of the state, substantially increasing the loss tolerance of the teleported information.\cite{ref51}

Furthermore, to combat the attenuation introduced during the initial distribution of the entangled resource state itself, the receiving node can pass the signal through a Noiseless Linear Amplifier (NLA) before final state reconstruction.\cite{ref52} The NLA probabilistically amplifies the state without adding the deterministic noise penalties required by classical linear amplification limits, significantly improving the overall fidelity of the teleported state and expanding the effective range of the P2P link.\cite{ref52}

\subsection{Multipartite States for Multi-User Mesh Routing}

Scaling this architecture from a single, isolated P2P link into a comprehensive metropolitan 6G backhaul mesh requires expanding beyond basic two-mode entanglement. Network architects leverage complex multipartite continuous-variable entangled states to enable dynamic, multi-user routing topologies.\cite{ref57}

\begin{itemize}
\item \textbf{Einstein-Podolsky-Rosen (EPR) States:} These provide the highest absolute fidelity for dedicated, low-error-rate P2P links.\cite{ref57}
\item \textbf{Greenberger-Horne-Zeilinger (GHZ) States:} When backhaul data must be multicast, GHZ-based protocols minimize the number of required entangled modes while preserving acceptable fidelity across multiple receiving nodes, optimizing limited quantum resources.\cite{ref57}
\item \textbf{Linear Cluster States:} For applications demanding highly controlled, dynamic routing across a grid, linear cluster states offer the optimal mathematical trade-off between resource overhead and fidelity.\cite{ref57} They supply the essential multi-party correlations required to seamlessly route massive 6G traffic through a complex, quantum-secured mesh network.\cite{ref19}
\end{itemize}

\section{Ecosystem Integration and the 6G Hardware Horizon}

The translation of these profound mathematical frameworks---from Python-simulated Lindbladian master equations and radiatively cooled PyTorch tensor models---into commercially deployable, hardware-ready P2P quantum backhaul systems requires the convergence of a highly specialized engineering ecosystem.\cite{ref6} The realization of the 6G quantum-enabled network relies heavily on cross-pollination between advanced telecommunications research institutes, integrated photonics foundries, and macroscopic hardware metrology.\cite{ref6}

At the vanguard of this ecosystem integration is the academic and corporate nexus driving 6G standards, heavily concentrated around entities like the UC San Diego Center for Wireless Communications (CWC) and Qualcomm.\cite{ref60} Advancing the deployment of sub-THz and massive MIMO platforms, researchers leverage photorealistic RF digital twins and AI-native 6G frameworks to accurately predict how highly directive beams will interact with unique physical environments prior to hardware deployment.\cite{ref15} These simulation testbeds are critical for optimizing the physical placement and topology of the P2P backhaul links, ensuring the neural-network-driven quantum metasurfaces are deployed in nodes that maximize Line-of-Sight (Rician) dominance while actively learning to mitigate Rayleigh scattering clutter.\cite{ref19}

Furthermore, physically constructing the ``Quantum ASIC'' holographic bus demands unprecedented advancements in materials packaging. The system requires consolidating hundreds of discrete optical elements, delicate 28nm Cryo-CMOS controllers, NbTi transmission lines, and massive 1,000-element lithium niobate meta-atom arrays into single, factory-aligned modules capable of operating in rugged outdoor conditions.\cite{ref6} This monumental engineering feat is being spearheaded by emerging integrated photonics startups---such as San Diego-based Monarch Quantum.\cite{ref64}

By applying advanced systems engineering and stringent Quality Management Systems (QMS), these entities are developing unified ``Quantum Light Engines''.\cite{ref66} They focus on co-packaged optics and photonic integrated circuits (PICs), fundamentally transforming sprawling, fragile laboratory laser benches into ruggedized, modular architectures.\cite{ref64} This scalable manufacturing ensures that the precision hardware necessary to generate, route, and measure TMSV states can be reliably deployed in the harsh terrestrial environments characteristic of modern cell tower backhaul locations, accelerating the time to quantum advantage.\cite{ref65}

\section{Conclusion}

The imperative to deliver hyper-scalable, unconditionally secure 6G connectivity relies fundamentally on the capabilities of the underlying Point-to-Point radio backhaul infrastructure. As exponential data demands exhaust traditional sub-6 GHz spectrums, the mandated migration to millimeter-wave and sub-terahertz frequencies exposes critical vulnerabilities at the physical layer, specifically catastrophic rain attenuation, atmospheric molecular absorption, and profound signal degradation induced by urban multipath Rayleigh fading. Traditional macroscopic antenna designs and classical heterodyne receivers are mathematically incapable of circumventing the strict thermal and environmental limitations of the terrestrial 300K environment without resorting to massive, power-intensive hardware configurations that violate the thermodynamic limits of the required cryogenic systems.

The definitive solution resides in a comprehensive, quantum-native physical layer transition. By embracing a Quantum ASIC paradigm driven entirely by Engineering-as-Code and Quantum Design Automation, networks can replace static, coaxial-heavy hardware with programmable, holographic metasurfaces managed by ultra-low dissipation 28nm FD-SOI Cryo-CMOS controllers. Leveraging Deep Neural Network surrogates trained via robust MLOps pipelines, these software-defined nodes map and optimize complex electromagnetic phase arrays in real-time. By utilizing narrow $\pi$-radian baselines, the system fundamentally respects the stringent sub-micro-watt power budgets necessary for cryogenic operation.

Crucially, implementing Microwave Quantum Illumination protocols completely redefines the physical limits of backhaul signal detection. By deterministically transmitting Two-Mode Squeezed Vacuum states and applying phase-conjugate joint measurements at the receiver, the network actively exploits resilient quantum discord to reject overwhelming multipath clutter, yielding massive, Heisenberg-limited gains in the Chernoff error exponent. Finally, by integrating passive radiative sky cooling via nanophotonic structures with dynamic 10 mK cryogenic overcoupling of the NbTi transmission lines, the architecture forcefully drains the terrestrial 300K thermal penalty from the communication channel. This extreme noise suppression provides the stable, quantum-coherent conduit necessary to execute Continuous-Variable Quantum Teleportation across distributed multi-user topologies using multipartite states. As integrated photonics ecosystems and advanced RF digital twins continue to mature, this synthesis of superconducting metamaterials, quantum optics, and thermodynamic engineering establishes the foundational blueprint for the future of global, quantum-secured terrestrial telecommunications.

\begin{thebibliography}{99}

\bibitem{ref1} Hybrid Satellite--Terrestrial Networks toward 6G: Key Technologies and Open Issues. MDPI. \url{https://www.mdpi.com/1424-8220/22/21/8544}. Accessed March 7, 2026.

\bibitem{ref2} Hybrid Satellite--Terrestrial Networks toward 6G: Key Technologies and Open Issues. PMC. \url{https://pmc.ncbi.nlm.nih.gov/articles/PMC9658377/}. Accessed March 7, 2026.

\bibitem{ref3} Wireless Terrestrial Backhaul for 6G Remote Access: Challenges and Low Power Solutions. Frontiers. \url{https://www.frontiersin.org/journals/communications-and-networks/articles/10.3389/frcmn.2021.710781/full}. Accessed March 7, 2026.

\bibitem{ref4} Millimetre-Wave Backhaul for 5G Networks: Challenges and Solutions. MDPI. \url{https://www.mdpi.com/1424-8220/16/6/892}. Accessed March 7, 2026.

\bibitem{ref5} The Benefits and Challenges of Microwave Backhaul. Infinity Technology Solutions. \url{https://infinitytdc.com/benefits-challenges-microwave-backhaul/}. Accessed March 7, 2026.

\bibitem{ref6} Synthesizing Quantum Metamaterial Paper.

\bibitem{ref7} Satellite-Based Continuous-Variable Quantum Communications: State-of-the-Art and a Predictive Outlook. University of Southampton. \url{https://www.southampton.ac.uk/~sxn/paper/cst-2018.pdf}. Accessed March 7, 2026.

\bibitem{ref8} A Tutorial on Non-Terrestrial Networks: Towards Global and Ubiquitous 6G Connectivity. \url{https://arxiv.org/html/2412.16611v1}. Accessed March 7, 2026.

\bibitem{ref9} Satellite-Terrestrial Quantum Networks and the Global Quantum Internet. arXiv.org. \url{https://arxiv.org/html/2410.13096v1}. Accessed March 7, 2026.

\bibitem{ref10} Performance Comparison of Rayleigh and Rician Fading Channels In QAM Modulation Scheme Using Simulink Environment. \url{https://www.ijceronline.com/papers/Vol3_issue5/Part.4/I035456062.pdf}. Accessed March 7, 2026.

\bibitem{ref11} Z. Elabdin, M. R. Islam, O. O. Khalifa. Mathematical model for the prediction of microwave signal attenuation due to duststorm. ResearchGate. \url{https://www.researchgate.net/profile/Md-Islam-101/publication/238074400_Mathematical_model_for_the_prediction_of_microwave_signal_attenuation_due_to_duststorm/links/55ac682108ae481aa7ff57cd/Mathematical-model-for-the-prediction-of-microwave-signal-attenuation-due-to-duststorm.pdf}. Accessed March 7, 2026.

\bibitem{ref12} A Review on Rain Signal Attenuation Modeling, Analysis and Validation Techniques: Advances, Challenges and Future Direction. MDPI. \url{https://www.mdpi.com/2071-1050/14/18/11744}. Accessed March 7, 2026.

\bibitem{ref13} CmWave and Sub-THz: Key Radio Enablers and Complementary Spectrum for 6G. Weizmann Institute of Science. \url{https://www.weizmann.ac.il/math/yonina/sites/math.yonina/files/CmWave\%20and\%20Sub-THz\%20Key\%20Radio\%20Enablers\%20and.pdf}. Accessed March 7, 2026.

\bibitem{ref14} Wireless Backhaul in 5G and Beyond: Issues, Challenges and Opportunities. Middle East Technical University. \url{https://open.metu.edu.tr/bitstream/handle/11511/101776/index.pdf}. Accessed March 7, 2026.

\bibitem{ref15} Unlocking Sub-THz by Robotic Aerial Base Stations: Joint Deployment and Wireless Backhaul Routing. IEEE Xplore. \url{https://ieeexplore.ieee.org/iel8/8782661/10362961/10766414.pdf}. Accessed March 7, 2026.

\bibitem{ref16} The Current Progress and Future Prospects of Path Loss Model for Terrestrial Radio Propagation. MDPI. \url{https://www.mdpi.com/2079-9292/12/24/4959}. Accessed March 7, 2026.

\bibitem{ref17} Microwave Terrestrial Link Rain Attenuation Prediction Parameter Analysis. Institute for Telecommunication Sciences. \url{https://its.ntia.gov/publications/download/84-148_ocr_reduced.pdf}. Accessed March 7, 2026.

\bibitem{ref18} Propagation Model Tuning for Terrestrial Microwave Links in Palestine. DSpace. \url{https://dspace.alquds.edu/bitstreams/a5ffdaf7-7263-4bf9-a454-066ece7d2f95/download}. Accessed March 7, 2026.

\bibitem{ref19} Performance Analysis of Rayleigh and Rician Fading Channel Models using Matlab Simulation. MECS Press. \url{https://www.mecs-press.org/ijisa/ijisa-v5-n9/IJISA-V5-N9-11.pdf}. Accessed March 7, 2026.

\bibitem{ref20} Performance Comparison of Rayleigh and Rician Fading Channel Models: A Review. IJAERD. \url{https://www.ijaerd.org/index.php/IJAERD/article/download/2355/2222}. Accessed March 7, 2026.

\bibitem{ref21} Unlocking the Power of Outdoor Wireless Networks: Strategies for Reliable and Scalable Connectivity. \url{https://www.aowireless.com/post/unlocking-the-power-of-outdoor-wireless-networks-strategies-for-reliable-and-scalable-connectivity}. Accessed March 7, 2026.

\bibitem{ref22} Multipath propagation characterization for terrestrial mobile and fixed microwave communication. Slideshare. \url{https://www.slideshare.net/slideshow/multipath-propagation-characterization-for-terrestrial-mobile-and-fixed-microwave-communication/66314119}. Accessed March 7, 2026.

\bibitem{ref23} Simulation of Multipath Fading Effects in Mobile Radio Systems. Microwave Journal. \url{https://www.microwavejournal.com/ext/resources/BGDownload/1/6/Multipath_Fading.pdf}. Accessed March 7, 2026.

\bibitem{ref24} Performance evaluation of Rayleigh and Rician Fading Channels using M-DPSK Modulation Scheme in Simulink Environment. IJERA. \url{https://www.ijera.com/papers/Vol3_issue3/HQ3313241330.pdf}. Accessed March 7, 2026.

\bibitem{ref25} Rayleigh fading. Wikipedia. \url{https://en.wikipedia.org/wiki/Rayleigh_fading}. Accessed March 7, 2026.

\bibitem{ref26} Capture Effect in the FSA-Based Networks under Rayleigh, Rician and Nakagami-m Fading Channels. MDPI. \url{https://www.mdpi.com/2076-3417/8/3/414}. Accessed March 7, 2026.

\bibitem{ref27} Introduction to Rayleigh Fading and Rician Fading. Medium. \url{https://medium.com/@acamvproducingstudio/introduction-to-rayleigh-fading-and-rician-fading-545e46e7a805}. Accessed March 7, 2026.

\bibitem{ref28} Rician Fading -- a Channel Model Often Misunderstood. Wireless Future Blog. \url{https://ma-mimo.ellintech.se/2020/03/02/rician-fading-a-channel-model-often-misunderstood/}. Accessed March 7, 2026.

\bibitem{ref29} Performance Analysis of Rayleigh and Rician Fading Channel Model. \url{https://gvpress.com/journals/IJMUE/vol12_no8/4.pdf}. Accessed March 7, 2026.

\bibitem{ref30} Spectrum issues on wireless backhaul. CIRCABC. \url{https://circabc.europa.eu/d/a/workspace/SpacesStore/cdd664eb-3007-4463-8dc9-e7940bd36d8f/RSPG15-607-Final_Report-Wireless_backhaul.pdf}. Accessed March 7, 2026.

\bibitem{ref31} Simulative analysis of Rayleigh and Rician fading channel model and its mitigation. Semantic Scholar. \url{https://www.semanticscholar.org/paper/Simulative-analysis-of-Rayleigh-and-Rician-fading-Kaur-Singh/e048f8970219c5bca7306cbaeda22661d79bd0e3}. Accessed March 7, 2026.

\bibitem{ref32} Mitigation of Fading Effects in Multipath Channels. Global Scientific Journal. \url{https://www.globalscientificjournal.com/researchpaper/MITIGATION-OF-FADING-EFFECTS-IN-MULTIPATH-CHANNELS.pdf}. Accessed March 7, 2026.

\bibitem{ref33} Quantum Microwave Illumination Protocols Enhancing Low-Power Radar Detection in Noisy Wireless Environments. Preprints.org. \url{https://www.preprints.org/manuscript/202601.2156/v1/download}. Accessed March 7, 2026.

\bibitem{ref34} Quantum Microwave Illumination Protocols Enhancing Low-Power Radar Detection in Noisy Wireless Environments. ResearchGate. \url{https://www.researchgate.net/publication/400217494_Quantum_Microwave_Illumination_Protocols_Enhancing_Low-Power_Radar_Detection_in_Noisy_Wireless_Environments}. Accessed March 7, 2026.

\bibitem{ref35} Point-to-Point Communication in Integrated Satellite-Aerial 6G Networks: State-of-the-Art and Future Challenges. IEEE Xplore. \url{https://ieeexplore.ieee.org/iel7/8782661/9309127/09466942.pdf}. Accessed March 7, 2026.

\bibitem{ref36} Quantum Microwave Illumination Protocols Enhancing Low-Power Radar Detection in Noisy Wireless Environments. Preprints.org. \url{https://www.preprints.org/manuscript/202601.2156}. Accessed March 7, 2026.

\bibitem{ref37} Microwave quantum illumination. PubMed. \url{https://pubmed.ncbi.nlm.nih.gov/25768743/}. Accessed March 7, 2026.

\bibitem{ref38} Microwave Quantum Illumination. DTIC. \url{https://apps.dtic.mil/sti/pdfs/AD1010145.pdf}. Accessed March 7, 2026.

\bibitem{ref39} Towards a Long-Range Microwave Quantum Radar. IEEE Aerospace and Electronic Systems Society. \url{https://ieee-aess.org/files/ieeeaess/slides/2025_IEEE\%20AESS\%20Distinguished\%20Lecture_Prof.\%20Livreri_19_08_2025.pdf}. Accessed March 7, 2026.

\bibitem{ref40} Microwave Quantum Illumination. DSpace@MIT. \url{https://dspace.mit.edu/handle/1721.1/95759}. Accessed March 7, 2026.

\bibitem{ref41} Radiative cooling to deep sub-freezing temperatures through a 24-h day--night cycle. PMC. \url{https://pmc.ncbi.nlm.nih.gov/articles/PMC5159822/}. Accessed March 7, 2026.

\bibitem{ref42} Realities of Space-Based Compute. Per Aspera. \url{https://www.peraspera.us/realities-of-space-based-compute/}. Accessed March 7, 2026.

\bibitem{ref43} Radiatively cooled quantum microwave amplifiers. ResearchGate. \url{https://www.researchgate.net/publication/382138149_Radiatively_cooled_quantum_microwave_amplifiers}. Accessed March 7, 2026.

\bibitem{ref44} Microwave Quantum Network Operates Resiliently up to 4 K. Bioengineer.org. \url{https://bioengineer.org/microwave-quantum-network-operates-resiliently-up-to-4-k/}. Accessed March 7, 2026.

\bibitem{ref45} Radiative sky cooling: fundamental physics, materials, structures, and applications. OSTI. \url{https://www.osti.gov/servlets/purl/1424949}. Accessed March 7, 2026.

\bibitem{ref46} Radiative cooling system: Topics by Science.gov. \url{https://www.science.gov/topicpages/r/radiative+cooling+system}. Accessed March 7, 2026.

\bibitem{ref47} Using outer space for cooling and energy on Earth. \url{https://iee.psu.edu/news/blog/using-outer-space-cooling-and-energy-earth}. Accessed March 7, 2026.

\bibitem{ref48} Radiative cooling to deep sub-freezing temperatures through a 24-h day--night cycle. ResearchGate. \url{https://www.researchgate.net/publication/311770331_Radiative_cooling_to_deep_sub-freezing_temperatures_through_a_24-h_day-night_cycle}. Accessed March 7, 2026.

\bibitem{ref49} Radiative sky cooling: Fundamental principles, materials, and applications. AIP Publishing. \url{https://pubs.aip.org/aip/apr/article/6/2/021306/570227/Radiative-sky-cooling-Fundamental-principles}. Accessed March 7, 2026.

\bibitem{ref50} Passive radiative cooling below ambient air temperature under direct sunlight. Stanford University. \url{https://web.stanford.edu/group/fan/publication/Raman_Nature_515_540_2014.pdf}. Accessed March 7, 2026.

\bibitem{ref51} Continuous variable quantum teleportation with sculptured and noisy non-Gaussian resources. arXiv. \url{https://arxiv.org/abs/0710.3259}. Accessed March 7, 2026.

\bibitem{ref52} Performance Analysis of Continuous Variable Quantum Teleportation with Noiseless Linear Amplifier in Seawater Channel. MDPI. \url{https://www.mdpi.com/2073-8994/14/5/997}. Accessed March 7, 2026.

\bibitem{ref53} Universal teleportation via continuous-variable graph states. \url{https://qnp.sjtu.edu.cn/userfiles/files/kxyj/Universal\%20teleportation\%20via\%20continuous-variable\%20graph\%20states.pdf}. Accessed March 7, 2026.

\bibitem{ref54} Quantum teleportation over continuous variables? Quantum Computing Stack Exchange. \url{https://quantumcomputing.stackexchange.com/questions/5187/quantum-teleportation-over-continuous-variables}. Accessed March 7, 2026.

\bibitem{ref55} Enhancing Continuous Variable Quantum Teleportation using Non-Gaussian Resources. arXiv. \url{https://arxiv.org/abs/2111.00672}. Accessed March 7, 2026.

\bibitem{ref56} Teleportation of hybrid entangled states with continuous-variable entanglement. PMC. \url{https://pmc.ncbi.nlm.nih.gov/articles/PMC9561706/}. Accessed March 7, 2026.

\bibitem{ref57} Analysis of multi-user quantum teleportation network based on continuous-variable entangled states. ResearchGate. \url{https://www.researchgate.net/publication/399489831_Analysis_of_multi-user_quantum_teleportation_network_based_on_continuous-variable_entangled_states}. Accessed March 7, 2026.

\bibitem{ref58} Advancing Standards. ETSI. \url{https://www.etsi.org/images/files/WorkProgramme/ETSI-Work-Programme-2024-2025.pdf}. Accessed March 7, 2026.

\bibitem{ref59} Quantum Technologies for Beyond 5G and 6G Networks: Applications, Opportunities, and Challenges. arXiv. \url{https://arxiv.org/html/2504.17133v1}. Accessed March 7, 2026.

\bibitem{ref60} The 2026 San Diego Wireless Summit. Jacobs School of Engineering. \url{https://jacobsschool.ucsd.edu/news/release/4060}. Accessed March 7, 2026.

\bibitem{ref61} Keysight and Qualcomm Advance RF Digital Twins to Scale Massive MIMO and AI-Native 6G Research. \url{https://www.keysight.com/us/en/about/newsroom/news-releases/2026/0303_pr26-039-keysight-and-qualcomm-advance-rf-digital-twins-to-scale-massive-mimo-and-ai-native-6g-research.html}. Accessed March 7, 2026.

\bibitem{ref62} UC San Diego Paper Defining Vision for 6G Integrated Sensing and Communication Among Most Downloaded. \url{https://today.ucsd.edu/story/uc-san-diego-paper-defining-vision-for-6g-integrated-sensing-and-communication-among-most-downloaded}. Accessed March 7, 2026.

\bibitem{ref63} Qualcomm accelerates 6G with AI-native device-to-data-center network transformation. \url{https://www.qualcomm.com/news/onq/2026/03/qualcomm-6g-device-to-data-center-transformation}. Accessed March 7, 2026.

\bibitem{ref64} Integrated Photonics Company Launches in San Diego. Monarch Quantum. \url{https://monarchquantum.com/2026/01/06/news20260106/}. Accessed March 7, 2026.

\bibitem{ref65} Monarch Quantum Selected to Deliver Integrated Photonics Quantum Light Engines$^{\mathrm{TM}}$ for NASA's First Planned Space-Based Quantum Gravity Gradiometer. PR Newswire. \url{https://www.prnewswire.com/news-releases/monarch-quantum-selected-to-deliver-integrated-photonics-quantum-light-engines-for-nasas-first-planned-space-based-quantum-gravity-gradiometer-302697215.html}. Accessed March 7, 2026.

\bibitem{ref66} Monarch Quantum Launches to Deliver Integrated Photonics Engines for the Quantum Ecosystem. \url{https://quantumcomputingreport.com/monarch-quantum-launches-to-deliver-integrated-photonics-engines-for-the-quantum-ecosystem/}. Accessed March 7, 2026.

\bibitem{ref67} Monarch Quantum: Home. \url{https://monarchquantum.com/}. Accessed March 7, 2026.

\bibitem{ref68} Australia's PIC market poised for rapid growth. Photonic Integrated Circuits News. \url{https://picmagazine.net/article/123239/Australia\%E2\%80\%99s_PIC_market_poised_for_rapid_growth}. Accessed March 7, 2026.

\end{thebibliography}

\end{document}
